Future work will focus on incorporating autonomous perception and adaptive 
control: integrating visual localization of parts, dynamic path planning
, and learning-based decision modules. By doing so, the SRRS framework can 
evolve from a manually orchestrated demonstration toward a truly self
-replicating robotic organism capable of operating in unstructured environments.

In conclusion, the SRRS project serves as both a scientific and engineering 
contribution: scientifically, as a platform for testing hypotheses about 
self-replicating mechanisms; and technically, as a fully functional ROS 
2–Gazebo system implementing attachment-based assembly. The methods and 
code we developed—spanning from the launch orchestration (robot.launch
) to the detailed node-level control (robotController.py, SRRScontrollerNode.
py, partController.py)-collectively provide a reproducible foundation for 
future research in autonomous manufacturing and self-replication robotics.
